% ===== tr-04 块1：第 4 章 pp.1--13；本片 pp.10--13（原文页 10--13）=====
% [原文 p.10]
于是它们诱导出一个定义良好的复结构，该复结构在 $\check x_j\mapsto\check x_j+1$ 下不变，从而在 $\check X(B)$ 上给出一个复结构。然后计算可得
\[
\omega=\sum d\check x_j\wedge dy_j=\frac{i}{2}\sum g^{jk}dz_j\wedge d\bar z_k,
\]
其中 $g_{ij}=\partial^2K/\partial y_j\partial y_k$。于是该度量是 Ricci 平坦的，当且仅当 $\det(g^{jk})=$ 常数，当且仅当 $\det(g_{jk})=$ 常数。\proofend
\end{proof}
% [Tutor Note: 此处原文 "where g_{ij} = ∂²K/∂y_j∂y_k" 左右指标不一致（左边 i,j，右边 j,k）；按后文 det(g_{jk}) 与度量分量定义，应为 g_{jk}=∂²K/∂y_j∂y_k。已 500 dpi 裁图核实原文确作 g_{ij}，正文照录，疑为原文笔误。]

照旧，我们把这个 Kähler 流形记作 $\check X(B,K)$。

这引出了如下定义。

\begin{definition}{4.2.4}
带有 Hessian 形式度量的仿射流形称为 \emph{Monge--Ampère 流形}，如果其局部势函数 $K$ 满足 Monge--Ampère 方程 $\det(\partial^2K/\partial y_i\partial y_j)=$ 常数。
\end{definition}

Hessian 流形与 Monge--Ampère 流形最早由 Cheng 与 Yau 在 [12] 中研究。

\begin{exercise}{4.2.5}
证明：由 Hessian 度量给出的 $\mathcal{T}_B$ 与 $\mathcal{T}_B^*$ 之间的等同，诱导出 Kähler 流形之间的典范同构 $X(B,K)\cong\check X(\check B,\check K)$，其中 $(\check B,\check K)$ 是 $(B,K)$ 的 Legendre 变换。
\end{exercise}

在镜像对称中出现一个关键的额外参数，称为 B-场（B-field）。它以非线性 sigma 模型中一个场的形式出现（该模型的靶空间为 Calabi--Yau 流形），在数学上是为了使镜像对称说得通所必需的。镜像对称粗略地断言 Calabi--Yau 流形 $X$ 的复模空间与 $\check X$ 的 Kähler 模空间之间的同构。若把 Kähler 模空间理解为 $\check X$ 上所有 Ricci 平坦 Kähler 形式构成的空间，则作为模空间只得到一个实流形，而我们需要一个复流形来与 $X$ 的复模空间相匹配。B-场被解释为一个元素 $\mathbf{B}\in H^2(\check X,\R/\Z)$，人们把 $\mathbf{B}+i\omega$ 看作 $\check X$ 上的复化 Kähler 类，其中 $\omega$ 是 $\check X$ 上的 Kähler 类。
% [Tutor Note: B-场是物理术语（B-field），对应的元素原文排作粗体（\mathbf{B}）；H^2(ˇX,R/Z) 只是复化 Kähler 类 B+iω 的「B 部分」，勿把 B 与 ω 混为同一个类。]

在我们这个玩具版本的镜像对称的语境中，我们把 B-场看作一个元素 $\mathbf{B}\in H^1(B,\Lambda_{\R}/\Lambda)$，其中 $\Lambda_{\R}=\Lambda\otimes_{\Z}\R$。这与上面关于 B-场的定义并不完全一致，因为这个群不一定与 $H^2(\check X,\R/\Z)$ 重合。然而，在许多重要情形，例如具有无挠整上同调的单连通 Calabi--Yau 三维流形，这两个群确实重合。更一般地，包括 K3 曲面与阿贝尔簇的情形，人们需要过渡到广义复结构 [43], [39], [7], [3], [44]，我们在此不想这样做。

注意到开集 $U$ 上 $\Lambda_{\R}/\Lambda$ 的一个截面可以被看作 $\mathcal{T}_U/\Lambda|_U$ 的一个截面，这样的截面通过平移作用在 $\mathcal{T}_U/\Lambda|_U$ 上，而这一作用关于半平坦复结构实际上是全纯的。于是，表示 $\mathbf{B}$ 的 Čech 1-上链 $(U_{ij},\beta_{ij})$ 允许我们通过诸交集 $U_{ij}$ 上的平移把 $X(B)$ 重新粘合起来。这是通过如下方式完成的：用由截面 $\beta_{ij}$ 给出的平移所定义的自同构来等同开子集 $f^{-1}(U_{ij})\subseteq f^{-1}(U_i)$ 与 $f^{-1}(U_{ij})\subseteq f^{-1}(U_j)$。
% [原文 p.11]
这给出一个新的复流形 $X(B,\mathbf{B})$。此外，若存在定义度量的多值势函数 $K$，则这些平移保持该度量，从而给出一个 Kähler 流形 $X(B,\mathbf{B},K)$。

于是镜像对称的完整玩具版本如下：

\begin{construction}{4.2.6}[玩具镜像对称构造]
设给定一个带有势 $K$ 的仿射流形 $B$（即对 $(B,K)$），以及 B-场 $\mathbf{B}\in H^1(B,\Lambda_{\R}/\Lambda)$、$\check{\mathbf{B}}\in H^1(B,\check\Lambda_{\R}/\check\Lambda)$。不难看出（如果你做过习题 2.5，你其实已经看到了），用 $B$ 上仿射结构定义的格局部系统 $\check\Lambda$，与用 $\check B$ 上仿射结构定义的格局部系统 $\Lambda$ 是同一个。于是我们说对
\[
(X(B,\mathbf{B},K),\check{\mathbf{B}})
\]
镜像于
\[
(X(\check B,\check{\mathbf{B}},\check K),\mathbf{B}).
\]
\end{construction}
% [Tutor Note: 构造 2.6 的「交叉」是关键：(X(B,B,K), ˇB) 镜像于 (X(ˇB,ˇB,ˇK), B)，左右两侧的 B-场互换；这正是 T 对偶在玩具模型中的体现，抄写时勿把两边的 B 与 ˇB 对应错。]

这在简单语境中给出了一个相当令人满意的镜像对称图景。在半平坦语境中还可以展开镜像对称的许多其他方面，见 [55] 与 [3] 第 6 章。这个半平坦情形是检验镜像对称中各种概念的理想试验场。然而归根结底，这只是对一般情形提供了有限的洞见。以这种方式出现的、带有半平坦 Ricci 平坦度量的紧 Calabi--Yau 流形只有复环面（由 Cheng 与 Yau 在 [12] 中证明）。为处理更有意思的情形，我们需要允许奇异纤维，从而允许 $B$ 的仿射结构出现奇点。奇异纤维的存在对于镜像对称中最有意思的方面是根本性的。

\section{带奇点的仿射流形}
为处理奇异纤维，我们引入如下定义。

\begin{definition}{4.3.1}
一个\emph{带奇点的（热带、整）仿射流形}是一个（$C^0$）流形 $B$，它有一个开子集 $B_0\subseteq B$，$B_0$ 上带有（热带、整）仿射结构，并且 $\Gamma:=B\setminus B_0$ 是余维数 $\ge2$ 的局部闭子流形的局部有限并。
\end{definition}

这里我们将给出这类带奇点仿射流形的一个相对简单的构造；更广的一类例子见 [22]；另见 [40] 与 [41]。

设 $\Delta$ 是 $M_{\R}=M\otimes_{\Z}\R$ 中的一个自反多面体，其中 $M=\Z^n$。这意味着 $\Delta$ 是一个格多面体，其内部有唯一的整点 $0\in\Delta$，并且极对偶多面体
\[
\nabla:=\{n\in N_{\R}\mid\langle m,n\rangle\ge-1\ \text{对所有}\ m\in\Delta\}
\]
也是格多面体。
% [Tutor Note: 自反多面体（reflexive polytope）的极对偶多面体记 ∇（nabla），与 §1--§2 中的对偶仿射流形 ˇB 无关；⟨m,n⟩ 是 M 与 N 的配对，m ∈ M_R, n ∈ N_R，方向勿反。]

% [原文 p.12]
设 $B=\partial\Delta$，并设 $\mathcal{P}$ 是 $B$ 的一个分解，其成员都是格多面体，即 $\mathcal{P}$ 是 $B$ 中格多面体构成的集合，满足：(1) $B=\bigcup_{\sigma\in\mathcal{P}}\sigma$；(2) $\sigma_1,\sigma_2\in\mathcal{P}$ 蕴含 $\sigma_1\cap\sigma_2$ 属于 $\mathcal{P}$，且是 $\sigma_1$ 与 $\sigma_2$ 的公共面；(3) 若 $\sigma\in\mathcal{P}$，则 $\sigma$ 的任一面属于 $\mathcal{P}$。

我们现在在 $B$ 上定义一个整仿射流形带奇点的结构，其判别轨迹 $\Gamma\subseteq B$ 定义如下。设 $\operatorname{Bar}(\mathcal{P})$ 表示 $\mathcal{P}$ 的第一重心重分，并设 $\Gamma\subseteq B$ 为 $\operatorname{Bar}(\mathcal{P})$ 中所有不含 $\mathcal{P}$ 的顶点（即零维胞腔）的单纯形、或与 $\mathcal{P}$ 的某个极大胞腔的内部相交的单纯形之并。置 $B_0:=B\setminus\Gamma$，我们如下在 $B_0$ 上定义仿射结构。$B_0$ 有开覆盖
\[
\{W_\sigma\mid\sigma\in\mathcal{P}\ \text{极大}\}\cup\{W_v\mid v\in\mathcal{P}\ \text{是顶点}\},
\]
其中 $W_\sigma=\operatorname{Int}(\sigma)$ 为 $\sigma$ 的内部，而
\[
W_v=\bigcup_{\substack{\tau\in\operatorname{Bar}(\mathcal{P})\\ v\in\tau}}\operatorname{Int}(\tau)
\]
是 $v$ 在 $\operatorname{Bar}(\mathcal{P})$ 中的（开）星。我们定义仿射坐标卡
\[
\psi_\sigma:W_\sigma\hookrightarrow\mathbb{A}_\sigma\subseteq N_{\R},
\]
它由 $W_\sigma$ 到 $\mathbb{A}_\sigma$ 的包含给出，这里 $\mathbb{A}_\sigma$ 表示 $N_{\R}$ 中包含 $\sigma$ 的唯一 $(n-1)$ 维仿射超平面。此外，取 $\psi_v:W_v\to M_{\R}/\R v$ 为投影。容易验证，对 $v\in\sigma$，$\psi_\sigma\circ\psi_v^{-1}$ 是整仿射线性的（整性来自 $\Delta$ 的自反性！），于是 $B$ 是一个整仿射流形带奇点。
% [Tutor Note: 判别轨迹 Γ 的定义是「Bar(P) 中不含 P 的零维胞腔（顶点）的单纯形」与「与 P 某极大胞腔内部相交的单纯形」的并，两个条件缺一不可；文字层把两条件连成一句，容易只抄到第一个。]

\begin{example}{4.3.2}
设 $\Delta\subseteq\R^4$ 为如下各点的凸包：
\begin{align*}
&(-1,-1,-1,-1),\qquad (4,-1,-1,-1),\\
&(-1,4,-1,-1),\qquad (-1,-1,4,-1),\qquad (-1,-1,-1,4).
\end{align*}
选取 $B=\partial\Delta$ 的一个剖分 $\mathcal{P}$ 为标准单纯形；这可以做得很有规律，使得 $\mathcal{P}$ 在每个二维面上的限制如图 4.1 中的浅色（细）线所示。这给出由图中的深色（粗）线所描绘的判别轨迹 $\Gamma$；从二维面边界伸出的线段意在说明包含在相邻二维面中的那部分判别轨迹。那里的判别轨迹并不包含在该二维面的平面内。特别地，就我们所定义的仿射结构而言，判别轨迹在 $\Gamma$ 位于 $\Xi$ 的棱上的顶点处不是平面的。注意 $\Gamma$ 是一个三价图，有两种三价顶点：刚刚提到的非平面顶点，以及位于二维面内部的平面顶点。
\end{example}
% [Tutor Note: 原文此句作 "the edges of Ξ"，但符号 Ξ 在本章此前从未出现；由上下文（B = ∂Δ）看极可能是 Δ 之误；500 dpi 裁图（p12 末段）确认原文字母确为 Ξ，正文照录，存疑待考。]

% [原文 p.13]
\figimg[0.72]{p4_fig01.png}
\figtitle{4.1}{}
% [Tutor Note: 原文图注仅有编号 "Figure 1."，没有描述文字，故图题第二参数留空；图中细线为剖分 P 的限制，粗线为判别轨迹 Γ。]

对一个仿射流形而言，格局部系统 $\Lambda$ 的单值（monodromy）是仿射结构的一个重要特征。在这个例子中，分析绕判别轨迹的环路所产生的单值非常有用。若 $v$ 是 $\Gamma$ 中位于 $\Delta$ 某个二维面内部的顶点，人们可以考虑在 $B_0$ 中基于 $v$ 附近的环路，它们围绕与 $v$ 相邻的三条 $\Gamma$ 的线段。计算这些单值矩阵在适当的基下具有如下形式，是一件令人愉快的练习：
\[
T_1=\begin{pmatrix}1&0&0\\1&1&0\\0&0&1\end{pmatrix},\quad
T_2=\begin{pmatrix}1&0&0\\0&1&0\\1&0&1\end{pmatrix},\quad
T_3=\begin{pmatrix}1&0&0\\-1&1&0\\-1&0&1\end{pmatrix}.
\]
它们是通过研究环路所经过的各坐标卡之间过渡映射的复合来计算的。这些矩阵可以被看作指定了把仿射结构延拓过 $v$ 在 $\Gamma$ 中某个邻域时的障碍。当然，$\check\Lambda$ 的单值是这些矩阵的转置逆。类似地，若 $v$ 是 $\Gamma$ 中位于 $\Delta$ 的一条棱上的顶点，则单值将取如下形式：
\[
T_1=\begin{pmatrix}1&-1&0\\0&1&0\\0&0&1\end{pmatrix},\quad
T_2=\begin{pmatrix}1&0&-1\\0&1&0\\0&0&1\end{pmatrix},\quad
T_3=\begin{pmatrix}1&1&1\\0&1&0\\0&0&1\end{pmatrix}.
\]
% [Tutor Note: 两组单值矩阵已在 600 dpi 下逐元素核对：第一组（二维面内部顶点）T_3 的第一列有负号 (−1,−1)，第二组（棱上顶点）T_3 的第一行为 (1,1,1) 且全为正；两组互为对偶关系的转置逆，符号极易抄错。]
于是我们看到，两类顶点的单值在 $\Lambda$ 与 $\check\Lambda$ 之间互换。

[20] 的一个主要结果是

\begin{theorem}{4.3.3}
若 $B$ 是一个三维带奇点的热带仿射流形，使得 $\Gamma$ 是三价的、且 $\Lambda$ 在每个顶点处的单值是上述两种类型之一，则 $f_0:X(B_0)\to B_0$ 可以紧化成拓扑纤维化 $f:X(B)\to B$。对偶地，$\check f_0:\check X(B_0)\to B_0$ 可以紧化成拓扑纤维化 $\check f:\check X(B)\to B$。$X(B)$ 与 $\check X(B)$ 都是拓扑流形。
\end{theorem}
% 注：定理 4.3.3 的最后一句（原文 "Both X(B) and ˇX(B) are topological manifolds."）在原文 p14 首行；为使定理陈述完整，随定理置于本片，p14 标记仍归于下一片 tr04d。